% Cover letter for the IEEE Access submission.
%
% Deliberately short. IEEE Access does not run a scope-fit negotiation the way
% the MDPI and Elsevier submissions did, and the editor reads the abstract
% anyway. What a cover letter can usefully carry here is exactly three things:
% what the paper claims, what makes it verifiable, and the disclosures that
% would otherwise surface as a query.
%
% No em dashes and no "we believe / we are confident" padding: the previous
% cover letters had both and neither survived a read-through.
\documentclass[11pt,a4paper]{article}
\usepackage[T1]{fontenc}
\usepackage[utf8]{inputenc}
\usepackage{mathptmx}
\usepackage[margin=2.4cm]{geometry}
\usepackage{microtype}
\pagestyle{empty}
% Date generated at build time so the letter cannot go stale between
% rebuilds. Redefined to day-month-year; the LaTeX default is US order.
\renewcommand{\today}{\number\day\ \ifcase\month\or January\or February\or March\or April\or May\or June\or July\or August\or September\or October\or November\or December\fi\ \number\year}
\setlength{\parindent}{0pt}
\setlength{\parskip}{7pt}

\begin{document}

{\raggedleft \today\par}

\vspace{4pt}
To the Editors, \emph{IEEE Access}

\vspace{6pt}
Dear Editors,

We resubmit the article \textbf{``Surrogate Selection for
Reinforcement-Learning HVAC Control: Step Size, Not Predictive Accuracy,
Predicts Control Utility''} (Manuscript ID Access-2026-42535) together with its
supplementary material.

\textbf{What was corrected.} The manuscript was returned to draft on 2 September
2026 for want of author biographies. A biography for each of the six authors is
now included, below the references. Reviewing the language as you suggested also
surfaced a separate defect: several abbreviations, including the two central to
the paper, were used before they were defined, and a few were never expanded.
Every abbreviation is now defined at its first use.

\textbf{What the article reports.} Reinforcement-learning controllers for
building HVAC systems are trained against fast learned surrogates, and the
surrogate is selected by predictive error. On the BOPTEST benchmark we find that
criterion to be not merely uninformative but inverted: the least accurate
surrogate in the study is the only one that trains a usable controller, while the
most accurate one trains a controller that leaves the comfort band for more than
77\% of the horizon on every random seed. We identify the property that does
predict training utility, the magnitude of the surrogate's per-step increment,
and isolate it with a rescaling control that holds response-surface roughness
fixed while changing only the integration step. We then give an architecture,
role-separated surrogate training, that acts on the finding, and a
receding-horizon MPC baseline showing that a planner reproduces the same inverted
ordering, so the effect belongs to optimizing through the surrogate rather than
to reinforcement learning specifically.

\textbf{Why it is verifiable.} Every experiment runs on BOPTEST, the open
IBPSA-derived benchmark, rather than on a private building model, so the results
can be reproduced and contested by other groups. Five hypotheses were written
down with their pass criteria before the corresponding runs; three are falsified,
including two of our own predictions, and all are reported as they came out. The
central comparisons are replicated across random seeds, and the censor-weight
sweep is replicated at every swept setting. One planned experiment was executed
and then invalidated by a defect we found in our own planner; that run is
recorded in the project artifacts rather than discarded, and the corrected run
supersedes it.

\textbf{Fit to IEEE Access.} The contribution is methodological, about how to
choose a training environment for a learning controller, demonstrated on a
building-control case study. It sits between machine learning and building
systems, and its central result is a negative one about a criterion the field
uses by default. IEEE Access states that it welcomes multidisciplinary and
applications-oriented work and negative results that do not fit the traditional
journals, which is the reason we bring the article here.

\textbf{Disclosures.} The work is original and is not under consideration
elsewhere. Earlier results from this study were presented as a conference talk at
the WCCM--ECCOMAS Congress 2026 (Munich, 19--24 July 2026); the present article
substantially extends that material and shares no text with it. All authors have
approved the submission and declare no competing interests. Author photographs
are not included and can be supplied if the article proceeds. The project
codebase is not publicly released at this stage; it
is available to the editors and reviewers on request, together with the derived
artifacts, from the corresponding author.

Thank you for considering the article.

\vspace{10pt}
Yours sincerely,

\vspace{6pt}
Almaz Sapargali, on behalf of all authors\\
Al-Farabi Kazakh National University and LLP Digit Alem, Almaty, Kazakhstan\\
almaz.sapargali2001@gmail.com

\end{document}
